% page 543 du fac-simile (image p-542 du scan)
% bloc 170.2 x 229.0 mm ; marge gauche 31.8 mm
\pgc{10.160mm}{0.000mm}{
\l{\cel{52}535}
\l{}
\l{\sou{skarabeo}. I. (\sou{zool.}) Nomo generala di ula insekti koleoptera,}
\l{\cel{3}partikulare ula lamelikornieri. - II. (\sou{geom.}) Kurvo qua pens-}
\l{\cel{3}igas pri skarabeo (insekto) - loko di la pedi di la perpendikli}
\l{\cel{3}trasita de punto ek la bisekanto di angulo orta til lineo kons-}
\l{\cel{3}tante-longa, di qua la extremaji quaze glitas sur la lateri di}
\l{\cel{3}la angulo. - EFIRS.}
\l{}
\l{\sou{skaramucho}. (\sou{teatro}). Rolo komika, bufonala, en la olima komedio}
\l{\cel{3}Italiana, di qua la vesto esis nigra. - DEFI.}
\l{}
\l{\sou{skarifikar}. (\sou{trans.}) (\sou{medicino}). Incizar surfacale (la pelo) per}
\l{\cel{3}lanceto, per bisturio, o per aparato ek kupra buxo di qua un}
\l{\cel{3}ek la edri esas provizita ye fendureti tra qui ekiras, per la}
\l{\cel{3}preso da kliko-resorto, plura lanceti qui esas muntita sur pi-}
\l{\cel{3}voto komuna.- EFIRS.}
\l{}
\l{\sou{skarlato}. Koloro dazlante reda, obtenita de la kochenilo di nopalo,}
\l{\cel{3}o de la kochenilo di querko, e qua uzesas por tintar. - DEFIRS.}
\l{}
\l{\sou{skarlatino}. (\sou{patol.}) Influro di la pelo, kontagiala, karakterizáta}
\l{\cel{3}da makuli skarlata. - DEFIRS.}
\l{}
\l{\sou{skarmuchar}. (\sou{trans. e netrans.)}\cel{2}(\sou{Lor milito}) (\sou{aludante detach}}
\l{\cel{3}\sou{mento, tiralieri di du armei}).Atakar per kombateti. - DEFIS.}
\l{}
\l{\sou{skarsa}.\cel{2}Qua ne abundas.\cel{2}EI.}
\l{}
\l{\sou{skeleto}. (\sou{anat.}) Osto karpenturo di animalo (vertebroza).- DEFIRS.}
\l{}
\l{\sou{skemo}. I. (a) Trasuro qua reprezentas, per la proporcioni e la}
\l{\cel{3}relati di ula figuri, la vario-legi di ula sorti de fenomeni -}
\l{\cel{3}en la domeni di fiziko, mekaniko, statistiko, e c.; (b) Trasuro}
\l{\cel{3}en qua la planeti diversa esas reprezentita en lia loko ye ta}
\l{\cel{3}o ca instanto ; (c) Trasuro qua reprezentas la dispozeso di}
\l{\cel{3}organo, di aparato, e c. - II. (\sou{yuro ekleziala)}. Propozajo,}
\l{\cel{3}redaktita formale, por submiseso a koncilo. - III. (\sou{filoz.})}
\l{\cel{3}Formo ye qua ni reprezentas,en nia mento, koncepton intelektala.}
\l{\cel{3}- DEFIRS.}
\l{}
\l{\sou{skeno.}\cel{1}Asembluro ek fili de kanabo, de silko, de lano, e c., vol-}
\l{\cel{3}vita ordinoze por ke li ne inter-intrikez. - E.}
\l{}
\l{skeptika. Qua profesas la dubito filozofiala, olqua konsistas ye}
\l{\cel{4}egardar raciono kom neapta, pro olua naturo ipsa, konocar}
\l{\cel{3}ulo certe. - DEFIRS.}
\l{}
\l{\cel{1}\sou{skerco}. (\sou{muziko)}.\cel{2}Muzikajo ne-longa, segun stilo jokema, petulema}
\l{\cel{4}lejera, lor la pleo di qua oportas, prefere, separar kam ligar}
\l{\cel{4}la noti. - DEFIR.}
\l{}
\l{skermar. (netrans.) Praktikar ta exerco qua kustumigas pri uzar}
}
